\section{Introduction}
\label{sec:intro}

% GOAL OF THIS SECTION: motivate the problem, state the gap, state what
% you built, and preview the contribution. ~1 page.

% --- The motivating tension (you already have strong wording for this
%     in _SHARED/SOP_master_draft.md — reuse and expand it here) ---
% A student handed the answer without doing the work learns less than
% one who struggled toward it. As AI tutors get better at producing
% correct answers, they may get worse for the students using them.
TODO: open with the learning-vs-answering tension.

% --- The gap in existing tools ---
% The problem is not access to AI tools, but the absence of guidance on
% using them without letting the tool do the thinking.
TODO: state the gap precisely.

% --- Your reframing ---
% The objective shifts from "minimize loss until the model is as capable
% as possible" to "constrain and guide the model toward making students
% more effective learners."
TODO: state the reframed objective.

% --- What you built (one paragraph, high level) ---
% A course-specific tutoring prototype that works through hundreds of
% pages of course material (slides, textbooks) and filters it down to the
% two or three references and images that actually answer the student's
% question, based on what was taught in that specific course --- pointing
% students back to the exact sources they were expected to study.
TODO: one-paragraph system overview.

% --- Contributions (bullet the 3-4 concrete things this report delivers) ---
\paragraph{Contributions.}
\begin{itemize}
  \item TODO: the problem framing / design principle.
  \item TODO: the system design.
  \item TODO: the implementation.
  \item TODO: the evaluation / findings (preliminary is fine).
\end{itemize}
